\section{Results}
\label{sec:results}

The four back-off constructions are compared across the disturbance sweep on the committed
twin data (seed $20260614$). Table~\ref{tab:regime} is the full regime map and
Figure~\ref{fig:violation} its headline projection: realised violation against disturbance
magnitude for every method, with the nominal target $\alpha=0.10$ marked. The reading is
consistent across the sweep, and we summarise it method by method below.

\subsection{The selection effect: marginal back-offs are unsafe under optimisation}
\label{sec:results-selection}

The normalised, adaptive back-off, the soft-sensor-style interval that is marginally valid by
construction, fails dramatically. Its realised violation sits between $0.43$ and $0.50$ across
the entire disturbance range against the $0.10$ target, roughly five times the nominal level,
with no dependence on $\sigma_z$. The failure is structural rather than a matter of the
disturbance being large. The mechanism is visible in the selected setpoint
(Figure~\ref{fig:fields}). At every $\sigma_z$ the adaptive method drives the optimiser to the
low-$D$ corner near $R\approx2.83$ and $D\approx33.7$, close to the deterministic optimum,
where the back-off is smallest relative to the true local risk. This is precisely the
operating point the selection-effect argument of Section~\ref{sec:methods-selection} predicts:
among feasible points the optimiser takes the most profitable, which is the one where the
marginally calibrated margin under-covers the conditional quantile. Making the interval
adaptive in width does not help, because the smooth scale model still leaves under-covered
pockets for the optimiser to occupy.

The constant margin is less catastrophic but still unsafe, and ultimately unusable. At small
disturbance it over-violates modestly, between $0.17$ and $0.19$ for $\sigma_z\le0.006$,
because a single margin sized to the pooled residual is too loose where the column is sensitive
and too tight where it is not. For $\sigma_z\ge0.015$ no constant margin admits a feasible
setpoint at all, since the margin required to cover the worst region of the box exceeds the
spec headroom everywhere. Neither marginal construction controls the violation rate that the
chance constraint demands.

\subsection{The oracle confirms the chance constraint is sound}
\label{sec:results-oracle}

That the marginal methods fail is a property of the back-off, not of the formulation. The
oracle back-off, the true conditional quantile, available here because the twin furnishes the
plant response, realises a violation between $0.079$ and $0.106$ across all seven disturbance
levels, that is approximately $\alpha$ by construction, and it remains feasible to the largest
$\sigma_z$ tested. The oracle is not a deployable method, since it needs the plant it is meant
to protect, but it certifies two things. The chance-constrained RTO of Eq.~\eqref{eq:chance} is
correctly posed, and a back-off that is conditionally correct controls the violation rate
exactly where the marginal ones do not. The oracle is therefore the bar the data-driven
conditional method must reach.

\subsection{Conditional calibration restores violation control}
\label{sec:results-fix}

The proposed method, CQR with a-posteriori calibration, tracks the oracle. Its realised
violation is between $0.063$ and $0.105$ for $\sigma_z\le0.020$, within about $0.01$ of the
oracle at every level and inside the target band, and it achieves this at near-oracle profit,
since the CQR and oracle profit columns of Table~\ref{tab:regime} differ by at most a few
dollars per hour at each $\sigma_z$. Two features distinguish it from the marginal baselines.
It selects setpoints that move with the disturbance, from $D\approx34.3$ at the realistic
$\sigma_z=0.006$ out to $D\approx36.6$ at $\sigma_z=0.020$, backing the column away from the
corner as the feed becomes more variable, exactly as the conditional risk requires and exactly
what the marginal methods fail to do. And where the residual finite-sample gap would otherwise
let it violate, the a-posteriori step closes it: at $\sigma_z=0.006$ the conditional quantile
alone is slightly optimistic, and the inflation $\kappa=1.53$ restores the target, with the
reported violation taken on an independent test draw. The conditional construction supplies the
right shape of back-off, and the a-posteriori step supplies the right scale.

\subsection{The boundary: conditional estimation runs out of data}
\label{sec:results-boundary}

The conditional method is not unconditionally safe. It degrades, gracefully and diagnosably, as
the disturbance outgrows the calibration set, and the a-posteriori inflation $\kappa$ is the
signal (Figure~\ref{fig:kappa}). For $\sigma_z\le0.010$ it stays near unity, between $1.0$ and
$1.5$, since the conditional estimate is already close to correct and little stretching is
needed. At $\sigma_z=0.015$ it jumps to $5.78$: the estimate must be stretched nearly sixfold to
be certified safe, and the method is operating at the edge of its feasible-$\kappa$ window. At
$\sigma_z=0.025$, more than four times the realistic disturbance, no feasible inflation certifies
the target and the method is reported infeasible, even though the oracle remains feasible there.
The gap between the two is the finite-sample cost of not knowing the plant. With the truth
surface the conditional quantile is exact; estimated from a calibration set whose support must
widen with $\sigma_z$, it eventually cannot be pinned down tightly enough for an optimiser to
lean on. The $\kappa$ trace makes this a readable property of the method rather than a silent
failure, since the operator sees the safety margin being stretched before it breaks.

\subsection{Safety, not profit}
\label{sec:results-profit}

A deliberate non-result is that the proposed method does not earn more money. At the realistic
disturbance the chance constraint is barely active: every method, including the deterministic
optimum at a profit of about \$5393\,h$^{-1}$, lies within $0.5\%$ of every other
(Table~\ref{tab:regime}; Figure~\ref{fig:profit}). The marginally higher profit of the unsafe
adaptive method is not an advantage. It is the profit of operating at a setpoint that violates
the spec half the time. The contribution is therefore not a profit improvement but calibrated
violation control at no cost in profit: the conditional method buys safety that the marginal
methods cannot deliver, while giving up essentially nothing in operating profit at the realistic
operating point. The profit separation between methods opens up only as $\sigma_z$ grows and the
constraint begins to bind, at which point the marginal methods are already either grossly
violating or infeasible.

\subsection{Surrogate validation and precision accounting}
\label{sec:results-valid}

Every violation rate above is read against the truth surface, so its credibility rests on that
surface interpolating the disturbance faithfully. The leave-one-feed-composition-slice-out test,
fitting on all $z\neq0.375$ and predicting the held-out $z=0.375$ slice, gives $R^2=0.917$ and a
mean absolute error of $0.0053$ in $x_B$ over $15$ held-out points (Figure~\ref{fig:validation}).
That error is the precision floor on the violation estimates. It is why the oracle lands between
$0.08$ and $0.11$ rather than exactly $0.10$, and it bounds how finely any method's violation can
be resolved. Layered on top is the Monte-Carlo sampling noise of the violation estimator itself,
of order $\pm0.01$ to $0.02$ across platforms; re-running the map on a different machine moves
individual rates by at most this much while leaving the regime structure and the fivefold
selection-effect gap intact. We therefore report violation rates as estimates with a band rather
than as point values. The qualitative claims, namely that the marginal back-off fails by about
fivefold, that the conditional method tracks the oracle, and that there is a data-starvation
boundary, all sit far outside that band.

\begin{figure}[t]
\centering
\includegraphics[width=\linewidth]{F3_violation_vs_sigma.png}
\caption{Realised constraint-violation rate versus disturbance magnitude $\sigma_z$ for the four
back-offs, with the nominal target $\alpha=0.10$ and the realistic $\sigma_z$ marked. Crosses
denote infeasible instances; shaded bands are the Monte-Carlo uncertainty. The marginal
back-offs (fixed, normalised) violate well above target, while CQR with a-posteriori calibration
tracks the oracle.}
\label{fig:violation}
\end{figure}

\begin{figure}[t]
\centering
\includegraphics[width=\linewidth]{F4_backoff_fields.png}
\caption{Back-off field $C(R,D)$ at the realistic disturbance for the naive-adaptive and the
CQR+a-posteriori methods, with each method's selected optimum ($\star$) and the oracle optimum.
The white curve is the true feasibility boundary. The adaptive margin stays small in the low-$D$
corner, so its optimum lands on the unsafe side; the CQR margin grows there and keeps its optimum
safe.}
\label{fig:fields}
\end{figure}

\begin{figure}[t]
\centering
\includegraphics[width=\linewidth]{F5_profit_vs_sigma.png}
\caption{Operating profit at the selected setpoint versus $\sigma_z$. At the realistic
disturbance all methods lie within half a percent of the deterministic optimum; the profit
separation opens only as the constraint begins to bind.}
\label{fig:profit}
\end{figure}

\begin{figure}[t]
\centering
\includegraphics[width=\linewidth]{F6_surrogate_validation.png}
\caption{Truth-surface validation. Left: predicted versus actual $x_B$ on the held-out
$z=0.375$ slice (leave-one-disturbance-out). Right: the monotone $x_B(R,D,z)$ response that
justifies the $z$-quantile oracle.}
\label{fig:validation}
\end{figure}

\begin{figure}[t]
\centering
\includegraphics[width=0.8\linewidth]{F7_kappa_vs_sigma.png}
\caption{A-posteriori inflation factor $\kappa^\star$ for the CQR method versus $\sigma_z$. The
margin is stretched (rising $\kappa$) before the method becomes infeasible, making $\kappa^\star$
a readable data-adequacy diagnostic.}
\label{fig:kappa}
\end{figure}

\begin{table}[t]
\centering
\caption{Regime map over disturbance magnitude $\sigma_z$ (canonical run, seed $20260614$,
leave-$z$-out $R^2=0.917$). For each method: operating profit (USD\,h$^{-1}$), realised violation
$\hat v$, a-posteriori inflation $\kappa$, and the selected setpoint $(R^\star,D^\star)$. INF
denotes an infeasible instance. The realistic disturbance is $\sigma_z=0.006$.}
\label{tab:regime}
\footnotesize
\begin{tabular}{llrrrrr}
\toprule
$\sigma_z$ & method & profit & $\hat v$ & $\kappa$ & $R^\star$ & $D^\star$ \\
\midrule
0.004 & oracle      & 5383 & 0.084 & 1.00 & 3.00 & 34.1 \\
      & CQR+apost   & 5381 & 0.063 & 1.00 & 3.00 & 34.1 \\
      & naive-fixed & 5388 & 0.190 & 1.00 & 3.00 & 33.9 \\
      & naive-adapt & 5394 & 0.429 & 1.00 & 2.83 & 33.7 \\
\midrule
0.006 & oracle      & 5374 & 0.079 & 1.00 & 3.00 & 34.3 \\
      & CQR+apost   & 5376 & 0.101 & 1.53 & 3.00 & 34.3 \\
      & naive-fixed & 5381 & 0.167 & 1.00 & 3.00 & 34.1 \\
      & naive-adapt & 5394 & 0.449 & 1.00 & 2.83 & 33.7 \\
\midrule
0.008 & oracle      & 5364 & 0.101 & 1.00 & 2.83 & 34.6 \\
      & CQR+apost   & 5364 & 0.097 & 1.00 & 2.83 & 34.6 \\
      & naive-fixed & 5374 & 0.159 & 1.00 & 3.00 & 34.3 \\
      & naive-adapt & 5396 & 0.472 & 1.00 & 3.00 & 33.7 \\
\midrule
0.010 & oracle      & 5351 & 0.106 & 1.00 & 2.80 & 34.9 \\
      & CQR+apost   & 5335 & 0.066 & 1.00 & 2.60 & 35.3 \\
      & naive-fixed & 5356 & 0.112 & 1.00 & 2.93 & 34.8 \\
      & naive-adapt & 5396 & 0.488 & 1.00 & 2.93 & 33.7 \\
\midrule
0.015 & oracle      & 5310 & 0.095 & 1.00 & 2.43 & 35.9 \\
      & CQR+apost   & 5298 & 0.105 & 5.78 & 3.00 & 36.2 \\
      & naive-fixed & \multicolumn{5}{c}{INF} \\
      & naive-adapt & 5396 & 0.480 & 1.00 & 3.00 & 33.7 \\
\midrule
0.020 & oracle      & 5288 & 0.087 & 1.00 & 2.20 & 36.5 \\
      & CQR+apost   & 5282 & 0.094 & 1.00 & 2.37 & 36.6 \\
      & naive-fixed & \multicolumn{5}{c}{INF} \\
      & naive-adapt & 5394 & 0.482 & 1.00 & 3.00 & 33.7 \\
\midrule
0.025 & oracle      & 5271 & 0.105 & 1.00 & 2.23 & 36.9 \\
      & CQR+apost   & \multicolumn{5}{c}{INF} \\
      & naive-fixed & \multicolumn{5}{c}{INF} \\
      & naive-adapt & 5396 & 0.502 & 1.00 & 3.00 & 33.7 \\
\bottomrule
\end{tabular}
\end{table}
